\chapter{Conclusão e Trabalhos Futuros}
\label{cap:conclusao}

Este capítulo consolida os principais resultados e conclusões obtidos ao longo do desenvolvimento do \textit{PID Simulator}, destacando as contribuições científicas, técnicas e pedagógicas alcançadas para a Engenharia de Computação e para o ensino de engenharia na Universidade Federal de Alagoas (UFAL). Por fim, discutem-se as limitações do trabalho e apresentam-se direcionamentos concretos para pesquisas e desenvolvimentos futuros.

\section{Considerações Finais}
\label{sec:consideracoes_finais}

O presente Trabalho de Conclusão de Curso propôs, implementou e validou experimentalmente uma plataforma web interativa para simulação física tridimensional, sintonia assistida e análise de telemetria de controladores PID em malha fechada. A motivação central esteve alicerçada na necessidade premente de superar os altos custos de licenciamento de ferramentas industriais consolidadas e mitigar a excessiva abstração matemática que permeia o ensino tradicional de sistemas de controle.

Ao longo do desenvolvimento, todos os objetivos específicos estabelecidos no Capítulo \ref{cap:introducao} foram plenamente atingidos:

\begin{enumerate}
    \item \textbf{Modelagem Física Realista:} Foram formuladas e implementadas equações dinâmicas e cinemáticas para três plantas distintas, incluindo veículo terrestre com perturbação topográfica de relevo 3D ($F_{\text{relevo}} = -mg\sin(\alpha)$), drone com sustentação vertical sob vento estocástico e navio sujeito a correntes marinhas oscilatórias;
    \item \textbf{Arquitetura Web de Alta Performance:} O frontend construído em Three.js e WebGL manteve taxa média estável de $59{,}4\text{ FPS}$, viabilizando uma experiência de instrumentação fluida sem a necessidade de instalação de qualquer software proprietário local;
    \item \textbf{Ambiente Seguro em Sandbox:} O editor de código acoplado ao construtor `new Function` demonstrou isolamento eficaz, permitindo a escrita de lógica de controle personalizada pelo estudante com interceptação de erros em tempo real;
    \item \textbf{Transparência Didática Multicanal:} O osciloscópio digital implementado possibilitou a visualização síncrona e a decomposição individual das forças proporcional ($u_P$), integral ($u_I$) e derivativa ($u_D$), permitindo ao discente correlacionar diretamente o papel de cada parcela com os fenômenos físicos observados na cena 3D (como a anulação do erro em rampas e a contenção do sobressinal);
    \item \textbf{Módulo Automatizado de Sintonia:} O assistente de sintonia em malha aberta automatizou com êxito a extração dos parâmetros FOPDT ($K, \tau, \theta$) e o cálculo de ganhos pelos critérios clássicos de Ziegler-Nichols e Cohen-Coon;
    \item \textbf{Infraestrutura Backend e Validação Numérica:} A API em FastAPI conteinerizada em Docker assegurou persistência estruturada, emissão rápida de relatórios técnicos em PDF com diagnósticos heurísticos e exportação de dados em CSV. A validação numérica comprovou erro quadrático médio de apenas $\text{RMSE} = 0{,}018\text{ m/s}$ frente à solução analítica contínua de referência.
\end{enumerate}

Dessa forma, comprova-se a hipótese de pesquisa formulada: a integração entre renderização tridimensional síncrona, decomposição explícita de sinais de controle e assistente de sintonia proporciona uma plataforma educacional superior, tangível e acessível para o aprendizado prático de engenharia de controle.

\section{Principais Contribuições}
\label{sec:principais_contribuicoes}

Dentre as principais contribuições técnico-científicas do trabalho destacam-se:

\begin{itemize}
    \item \textbf{Democratização do Acesso Laboratorial:} Disponibilização de um laboratório virtual livre, leve e aberto (\textit{open source}), capaz de ser executado em qualquer computador ou dispositivo com navegador web, viabilizando o ensino remoto e práticas extraclasse na UFAL;
    \item \textbf{Inovação na Decomposição Pedagógica de Sinais:} Abordagem pioneira que expõe visualmente cada termo ($u_P, u_I, u_D$) em canais de osciloscópio independentes com limites de saturação, transformando o conceito abstrato de esforço de controle em grandezas mensuráveis;
    \item \textbf{Convergência entre Controle e Engenharia de Computação:} Demonstração prática de como técnicas modernas de computação gráfica na web, conteinerização e microsserviços assíncronos podem ser aplicadas para resolver desafios pedagógicos de engenharias tradicionais.
\end{itemize}

\section{Limitações Identificadas}
\label{sec:limitacoes}

A despeito dos resultados plenamente satisfatórios, identificaram-se limitações operacionais que definem as fronteiras do presente estudo:

\begin{enumerate}
    \item \textbf{Simplificação de Modelos de Atrito:} Adotou-se prioritariamente o atrito viscoso linear ($b \cdot v$), omitindo-se não-linearidades complexas como o atrito de Coulomb estático com zona morta e o efeito Stribeck;
    \item \textbf{Execução na Thread Principal:} A execução da simulação e do código do usuário no escopo principal do navegador compartilha processamento com o ciclo de redesenho do DOM, o que pode gerar pequenas quedas de quadros caso o estudante introduza acidentalmente algoritmos computacionalmente pesados;
    \item \textbf{Inexistência de Acoplamento com Hardware Físico:} A plataforma restringe-se atualmente ao domínio da simulação puramente virtual, sem interface direta para controle de atuadores físicos reais em bancada.
\end{enumerate}

\section{Trabalhos Futuros}
\label{sec:trabalhos_futuros}

Como perspectivas de continuidade e expansão desta pesquisa, propõem-se os seguintes desdobramentos:

\begin{itemize}
    \item \textbf{Migração para Web Workers:} Isolar a execução do motor de física e da \textit{sandbox} de código em \textit{Web Workers} dedicados rodando em \textit{threads} secundárias em segundo plano, liberando integralmente a \textit{thread} principal para a renderização gráfica via `OffscreenCanvas`;
    \item \textbf{Implementação de Algoritmos Inteligentes de Sintonia:} Incorporar módulos de sintonia automática baseados em meta-heurísticas de inteligência artificial, tais como Algoritmos Genéticos (GA), Otimização por Enxame de Partículas (PSO) e Aprendizado por Reforço (\textit{Reinforcement Learning});
    \item \textbf{Integração com Hardware Real (\textit{Hardware-in-the-Loop}):} Utilizar a API WebSerial ou WebSockets para permitir que o simulador envie o esforço de controle e receba a telemetria de microcontroladores reais (como ESP32 ou Arduino) conectados à porta USB do computador;
    \item \textbf{Modelagem de Novas Plantas Clássicas:} Expandir a biblioteca de experimentos para incluir plantas multivariáveis ou instáveis em malha aberta consagradas no ensino, tais como o pêndulo invertido sobre carro móvel, o levitador magnético e sistemas de múltiplos tanques acoplados;
    \item \textbf{Avaliação Pedagógica em Sala de Aula:} Conduzir um estudo de usabilidade empírica com discentes de graduação da UFAL, quantificando o ganho de aprendizado através de questionários pré e pós-teste em turmas regulares de Sistemas de Controle.
\end{itemize}
